\section{Introduction}
\label{sec:introduction}

% (~4-5 pages)

% Opening: The gut microbiome as a complex adaptive system
% - 100 trillion microorganisms, 1000+ species
% - Not just passengers — active participants in host physiology
% - Metabolic, immune, and neurological functions

% The dysbiosis problem
% - Definition: disruption of the symbiotic equilibrium
% - Clinical relevance: IBD, obesity, diabetes, neurological disorders
% - The cascade: dysbiosis → SCFA depletion → barrier degradation →
%   endotoxin translocation → systemic inflammation → neuroinflammation

% Limitations of current approaches
% - Alpha/beta diversity metrics capture "how much" but not "what structure"
% - Differential abundance (DESeq2) identifies individual taxa but misses
%   community-level organization
% - Co-occurrence networks capture pairwise relationships but not
%   higher-order topology

% The TDA opportunity
% - Persistent homology captures multi-scale topological features
% - Proven in other complex systems: financial markets, protein structure,
%   neural networks, materials science
% - Key parallel: correlation homogeneity → topological stability
%   (markets ↔ microbial ecosystems)

% The gut-brain axis connection
% - Vagal nerve as bidirectional communication highway
% - Neuropeptides produced/modulated by gut microbiota:
%   serotonin (95% produced in gut), ghrelin, CCK, VIP
% - Intestinal permeability as the gateway mechanism
% - Environmental triggers: mycotoxins, antibiotics, diet

% Our contribution
% - First application of persistent homology to detect regime shifts
%   in gut microbial co-occurrence networks
% - Topological biomarkers that correlate with permeability markers
%   and neuropeptide levels
% - Sliding-window regime detection as early warning system
% - Mycotoxin-induced dysbiosis as case study (CIRS connection)

% Paper outline
% Section 2: Background on microbiome, TDA, gut-brain axis
% Section 3: Theoretical framework linking topology to ecology
% Section 4: Methods (datasets, pipeline, statistics)
% Section 5: Results
% Section 6: Discussion (interpretation, CIRS implications, limitations)
% Section 7: Conclusion
